% Generated draft: requires mathematical and source review.
\documentclass[11pt]{article}
\usepackage[margin=25mm]{geometry}
\usepackage{fontspec}
\setmainfont{Latin Modern Roman}
\usepackage{amsmath,amssymb,mathtools,unicode-math}
\setmathfont{Latin Modern Math}
\usepackage{enumitem,needspace,xurl,hyperref}
\setlength{\parindent}{0pt}
\setlength{\parskip}{5pt}
\setlength{\emergencystretch}{3em}
\newcommand{\N}{\mathbb{N}}
\newcommand{\Z}{\mathbb{Z}}
\newcommand{\Q}{\mathbb{Q}}
\newcommand{\R}{\mathbb{R}}
\newcommand{\C}{\mathbb{C}}
\newcommand{\F}{\mathbb{F}}
\newcommand{\A}{\mathbb{A}}
\newcommand{\PP}{\mathbb{P}}
\newcommand{\E}{\mathbb{E}}
\newcommand{\eps}{\varepsilon}
\newcommand{\dd}{\,\mathrm{d}}
\begin{document}
\section*{UnsolvedMath ID: 30000413}
\textbf{Source problem code:} catalogue entry
\subsection{Definitions and mathematical statement}
Let $K/k$ and $L/l$ be regular function fields over the ``reasonable'' base fields intended by the report. A standard function-field convention makes these finitely generated extensions; regularity means that $k$ is relatively algebraically closed in $K$ and $K/k$ is separably generated, and likewise for $L/l$. The intended transcendence-degree range and the class of reasonable bases are not supplied.

For fields $F_1,F_2$, write $F_1\equiv F_2$ when they satisfy the same first-order sentences in the language of rings $\{0,1,+,-,\cdot\}$. This is the usual pure-field interpretation; the source does not specify an expanded language or predicates naming the base fields.

Assuming $K\equiv L$, the source asks both:
\begin{enumerate}
\item Must $k\equiv l$?
\item If in addition $k\cong l$, must $K\cong L$ as fields?
\end{enumerate}
Equivalently, with all function-field and base-class hypotheses retained, the implications sought are
\[
 K\equiv L\quad\Longrightarrow\quad k\equiv l,
\]
\[
 (K\equiv L)\ \text{and}\ (k\cong l)
       \quad\Longrightarrow\quad K\cong L.
\]
The elementary-equivalence assumption is inherited by the second question; it is not dropped merely because its edited version starts a new sentence. The field isomorphisms are abstract isomorphisms as stated, not a demand that an arbitrarily prescribed base isomorphism extend to $K\to L$.

Without a definition of reasonable bases, these are conditional questions with an unresolved domain of quantification, not assertions about every field. No arithmetic, geometric, characteristic, perfectness, or transcendence-degree restriction is silently selected. If the intended language instead names the extensions or their base fields, that would materially alter the question and must be recovered from the report.
\subsection{Short English statement}
Assuming the two regular function fields have the same first-order field theory, the question asks whether their bases do too, and whether isomorphic bases then force the function fields themselves to be isomorphic. The undefined reasonable-base class and language conventions remain essential.
\subsection{Sources}
\begin{itemize}
\item Local UnsolvedMath registry, permanent ID 30000413, code catalogue entry, ``Elementary Equivalence of Regular Function Fields.'' Proposed year 2006; proposer unspecified.
\item \emph{The Arithmetic of Fields} (2006), Oberwolfach Reports, pp.\ 317--378, DOI 10.4171/owr/2006/06; \url{https://doi.org/10.4171/owr/2006/06}.
\item Supplied reference labeled ``Pop primary paper'': \url{https://www2.math.upenn.edu/~pop/Research/files-Res/AWS03-fin.pdf}.
\item Statement status corrected\_verified; the supplied verification describes meaning-preserving editorial and LaTeX cleanup without added source-dependent content.
\item Literature triage checked 2026-08-21: partially solved. It names Pop, Rumely, Poonen, and related work as settling substantial arithmetic and geometric cases, while saying the general reasonable-base statement is broader. No precise theorems or separate references for Rumely and Poonen are supplied. Created 2026-08-25; updated 2026-09-16. Status and cited claims require independent verification.
\item Registry curation by Ulam AI / UnsolvedMath Contributors, CC BY 4.0; underlying report and literature provenance remain as supplied. Unreviewed statement-only draft.
\end{itemize}
\end{document}
